\section{Chỉ dẫn cell ba notebook đối chứng}
SVM, Logistic Regression và LSTM mỗi file có 48 cell, gồm 31 cell mã đã thực thi. Số cell dưới đây tính từ 1 và bao gồm markdown. Cấu trúc chung thuận tiện đối chiếu nhưng mỗi notebook tự chứa mã, cache, gói mô hình và output riêng; không đọc mô hình của notebook khác.
\begin{table}[H]\centering\small\caption{Các nhóm cell của ba notebook đối chứng.}
\begin{tabular}{r p{112mm}}\toprule Cell & Nội dung và bằng chứng\\\midrule
3 & Cấu hình mô hình, seed, split và chế độ fresh\\
5--8 & Thư viện, helper, tiến độ và đường dẫn\\
10--11 & Chính sách hệ thống chuyên mục và xác minh dataset sạch\\
13--15 & Thiết bị, đọc 20 Parquet, kiểm tra URL\\
17--21 & Phân bố, độ dài và thống kê dữ liệu đầu vào\\
23--24 & Chuẩn hóa PyVi, loại bỏ bài trùng lặp và ví dụ văn bản\\
26--27 & Chia 70/15/15, số mẫu, overlap và độ dài\\
29--30 & SVM/LR: từ dừng, TF--IDF; BiLSTM: từ vựng, packed sequence\\
32 & Tìm cấu hình trên validation và giữ mô hình/bản lưu mô hình chọn\\
34 & Fit temperature chỉ trên validation\\
35--36 & Score test, chỉ số, bootstrap, ma trận và dự đoán từng bài\\
38--39 & Xuất gói suy luận, nạp lại và đối chiếu trên ba bài validation\\
41--43 & Hiệu chỉnh xác suất test, lớp yếu và lỗi độ tin cậy cao\\
44 & SVM/LR: đặc trưng; BiLSTM: lịch sử, OOV và tỷ lệ từ giữ\\
46--47 & Hình đánh giá, PR từng lớp, ROC, manifest và bảng cấu hình chọn\\\bottomrule
\end{tabular}\end{table}

\section{Đường dẫn bằng chứng so sánh}
Các file \texttt{result/<model>/evaluation/metrics.json} và \texttt{comparison\_summary.json} lưu cấu hình chọn, mã nhận diện dữ liệu, chỉ số và phép đo thời gian. \texttt{split\_manifest.json} lưu chỉ số ba tập; \texttt{test\_scores.npz} và \texttt{predictions.parquet} cho phép tái tính chỉ số mà không chạy mô hình. \texttt{classification\_report.json} và \texttt{confusion\_matrix.json} giữ số liệu theo lớp. Tên \texttt{<model>} là \texttt{svm}, \texttt{logistic} hoặc \texttt{lstm}; PhoBERT dùng \texttt{result/evaluation/}.

Các bảng so sánh và khoảng ghép cặp được tính từ cùng thứ tự test, đủ 20 nhãn, theo phương pháp ở mục~\ref{sec:paired_comparison}. Hình báo cáo đối chứng là bản sao nguyên byte từ \texttt{result/<model>/figures/}; hash được đối chiếu với \texttt{figure\_manifest.json}. Source LaTeX giữ các bảng kết quả và caption cùng đơn vị đo, không thay output của notebook bằng một giá trị suy đoán.

\section{Ma trận số lượng của các đối chứng}
Ma trận dưới đây dùng số bài thay cho phần trăm. Tổng mỗi ma trận bằng 29,995, tổng hàng khớp số mẫu trong bảng chia tập. Các ma trận chuẩn hóa của cùng mô hình được đặt ở chương~\ref{chap:comparison}. Mỗi ma trận là một PNG riêng của notebook, dùng cùng ID 0--19 với PhoBERT.
\PortraitFigure{assets/baselines/svm/evaluation/confusion_matrix.png}{Ma trận số lượng của SVM; cùng test và hệ thống chuyên mục 20 lớp.}
\PortraitFigure{assets/baselines/logistic/evaluation/confusion_matrix.png}{Ma trận số lượng của Logistic Regression; tổng 29,995 bài.}
\PortraitFigure{assets/baselines/lstm/evaluation/confusion_matrix.png}{Ma trận số lượng của BiLSTM; mỗi hàng là một lớp tham chiếu.}
